\clearpage\pagestyle{otMain}\setlength{\otBlocEspacement}{2.5mm}\fontsize{11}{15}\selectfont\corrigesChapitre{9}{Programmes de calcul et dépendance}
\begin{blocLivre}{cor-9-1}\begin{corrige}{1}{Appliquer un programme}\demarche{Effectuer les étapes dans l’ordre.}
Pour 4 : $4\times3-5=\res7$. Pour $-2$ : $(-2)\times3-5=-6-5=\res{-11}$. Pour 1,5 : $4{,}5-5=\res{-0{,}5}$.
\end{corrige}\end{blocLivre}
\begin{blocLivre}{cor-9-2}\begin{corrige}{2}{Du programme à l’expression}\demarche{Placer des parenthèses autour du résultat intermédiaire.}
On ajoute 6 : $x+6$. On multiplie le résultat entier par $-2$ : $\res{-2(x+6)}$. En développant : $\res{-2x-12}$. Les parenthèses traduisent « le résultat ».
\end{corrige}\end{blocLivre}
\begin{blocLivre}{cor-9-3}\begin{corrige}{3}{Remonter un programme}\demarche{Inverser les opérations dans l’ordre inverse.}
En remontant : $35:5=7$, puis $7+4=\res{11}$. Avec une équation : $5(x-4)=35$ donne $x-4=7$, donc $x=11$. Contrôle : $(11-4)\times5=35$.
\end{corrige}\end{blocLivre}
\begin{blocLivre}{cor-9-4}\begin{corrige}{4}{Une formule de tarif}\demarche{Additionner partie fixe et partie variable.}
La partie fixe vaut 5 et la partie liée au temps $3t$ : $\res{P(t)=5+3t}$. On remplace $t$ : $P(0)=\res5$, $P(2)=5+6=\res{11}$ et $P(4)=5+12=\res{17}$ euros.
\end{corrige}\end{blocLivre}
\begin{blocLivre}{cor-9-5}\begin{corrige}{5}{Lire un couple}\demarche{Suivre la verticale puis l’horizontale.}
On part de 2 sur l’axe du temps, on rejoint le tracé puis on lit l’axe vertical : $\res{T(2)=60\,^\circ\mathrm C}$. Les unités viennent des axes.
\label{fig-cor-9-5}\figureGeo[corrige]{four}
\end{corrige}\end{blocLivre}
\begin{blocLivre}{cor-9-6}\begin{corrige}{6}{Lire une durée}\demarche{Partir cette fois de l’axe vertical.}
L’horizontale de niveau 100 coupe le tracé à l’abscisse $\res4$ min pour la première fois. On lit cette fois un temps, et non une température.
\label{fig-cor-9-6}\figureGeo[corrige]{four}
\end{corrige}\end{blocLivre}
\begin{blocLivre}{cor-9-7}\begin{corrige}{7}{Une lecture inverse double}\demarche{Parcourir toute l’horizontale.}
L’horizontale de niveau 60 coupe le tracé deux fois : à $\res2$ min pendant la chauffe et à $\res{10}$ min pendant le refroidissement. Il faut chercher toutes les intersections dans l’intervalle représenté.
\label{fig-cor-9-7}\figureGeo[corrige]{four}
\end{corrige}\end{blocLivre}
\begin{blocLivre}{cor-9-8}\begin{corrige}{8}{Reconnaître une relation}\demarche{Vérifier le passage par l’origine.}
La droite représentée passe par le point $(0\,;2)$, pas par l’origine. La relation n’est donc $\res{\text{pas proportionnelle}}$, même si les points sont alignés : pour zéro unité horizontale, la valeur verticale vaut déjà 2.
\label{fig-cor-9-8}\figureGeo[corrige]{graph-non-proportion}
\end{corrige}\end{blocLivre}
\begin{blocLivre}{cor-9-9}\begin{corrige}{9}{Construire un graphique}\demarche{Placer les couples avant de relier.}
Les prix sont $\res{0,\ 3,\ 6,\ 9}$ euros. Les points $(0\,;0)$, $(2\,;3)$, $(4\,;6)$ et $(6\,;9)$ sont alignés avec l’origine : c’est une proportionnalité. La corde pouvant être achetée en longueur intermédiaire, on trace le segment correspondant.
\label{fig-cor-9-9}\figureGeo[corrige]{graph-corde}
\end{corrige}\end{blocLivre}
\begin{blocLivre}{cor-9-10}\begin{corrige}{10}{Aire d’un carré}\demarche{Comparer l’effet d’un doublement.}
$A(c)=c^2$. Les aires valent $\res{1,\ 4,\ 9,\ 16}$ cm$^2$. Doubler le côté de 1 à 2 multiplie l’aire par 4 et non par 2 : l’aire $\res{\text{n’est pas proportionnelle}}$ au côté. Les points ne sont pas alignés.
\label{fig-cor-9-10}\figureGeo[corrige]{graph-carre}
\end{corrige}\end{blocLivre}
\begin{blocLivre}{cor-9-11}\begin{corrige}{11}{Deux programmes égaux}\demarche{Écrire les deux expressions avec la même lettre.}
A donne $4(x+3)-12=4x+12-12=4x$, comme B : les programmes sont $\res{\text{équivalents}}$. Pour obtenir 26, on résout $4x=26$ : $\res{x=6{,}5}$.
\end{corrige}\end{blocLivre}
\begin{blocLivre}{cor-9-12}\begin{corrige}{12}{Un tableur}\demarche{Remplacer la lettre par la cellule de la ligne.}
On saisit $\res{\texttt{=5+3*A2}}$. Lors de la recopie, A2 devient A3, A4, etc. : B3 affiche $\res{11}$ et B4 affiche $\res{17}$. La référence relative suit la durée de chaque ligne ; les nombres 5 et 3 restent fixes.
\label{fig-cor-9-12}\figureGeo[corrige]{tableur-prix}
\end{corrige}\end{blocLivre}
\begin{blocLivre}{cor-9-13}\begin{corrige}{13}{Retirer des carrés}\demarche{Soustraire les quatre aires retirées.}
L’aire initiale vaut 180 cm$^2$. Les quatre carrés retirent $4x^2$, donc $\res{A(x)=180-4x^2}$. Pour $x=2$, $A(2)=180-16=\res{164\ \mathrm{cm}^2}$. La condition $x<5$ évite que deux carrés d’un même petit côté se chevauchent.
\label{fig-cor-9-13}\figureGeo[corrige]{coins-retires}
\end{corrige}\end{blocLivre}
\begin{blocLivre}{cor-9-14}\begin{corrige}{14}{Deux courbes de prix}\demarche{Confirmer la lecture par le calcul.}
Les tracés se coupent au point $(2\,;8)$. L’égalité $4+2t=4t$ donne $4=2t$, donc $\res{t=2}$ h et un prix de $\res8$ euros. Avant 2 h, B est moins cher ; après 2 h, A est moins cher sur l’intervalle étudié.
\label{fig-cor-9-14}\figureGeo[corrige]{graph-tarifs}
\end{corrige}\end{blocLivre}
\begin{blocLivre}{cor-9-15}\begin{corrige}{15}{Une réserve qui diminue}\demarche{Associer une formule à un domaine réaliste.}
Le volume restant est $\res{V(t)=120-8t}$ litres. Il s’annule pour $120-8t=0$, soit $t=15$ min. Le modèle de vidange est valable pour $\res{0\le t\le15}$ ; un volume négatif ne représente pas la situation.
\label{fig-cor-9-15}\figureGeo[corrige]{graph-reserve}
\end{corrige}\end{blocLivre}
